\scanpage{47}
\uline{\unsure{Conséquence|Conséquences}}: Si $\bigcup_{A\in\Sigma}A^*=S$, alors \deleted{les} $(A^*)_{A\in\Sigma}$ est une partition \deleted{\unsure{des|sans|dans}} \unsure{non} panachée de $S$. Si $\bigcup_{A\in\Sigma}A^*\ne S$, alors posant
\[
 \deleted{\unsure{R,\Sigma|R\ \Sigma}}\mathpunct{.}\qquad
 \left\{\begin{aligned}
 R&=S\setminus\bigcup_{A\in\Sigma}A^*,\\
 R^*&=R\setminus\bigcup_{\substack{A\in\widehat\Sigma\\ A\subsetneq R}}A=R
 \end{aligned}\right.\qquad
 \begin{aligned}
 \widehat\Sigma&=\Sigma\cup\{R\}\\
 &\deleted{\unsure{\Sigma|\in\Sigma|\ne\Sigma_0}}
 \end{aligned}
\]
alors $(A^*)_{A\in\widehat\Sigma}$ est une partition non panachée de $S$.

\uline{NB} Si $R$ est lui-= un segment, alors \deleted{\unsure{l'intérêt} \unsure{de le prendre|de le joindre}} $\widehat\Sigma$ satisfait les conditions de $\Sigma$ -- et $R$ est un élément \deleted{``\unsure{sp.|spéc.}} ``isolé'' de $\widehat\Sigma$, en ce sens que \inserted{pour} \unsure{tt|$A$|$\forall$} autre \unsure{élément|élém} \inserted{$A$} de $\widehat\Sigma$, \deleted{\unsure{on}}. $R\cap A=\phi$ -- \deleted{\unsure{donc}} \inserted{i.e.} $R$ est \unsure{:|à} la fois minimal, et maximal. Considérons
\[
 \begin{aligned}
 (\sigma\ 3)\qquad R=S\setminus\bigcup_{A\in\Sigma}A\;
 \bigl(=S\setminus\bigcup A^*\bigr)&\quad\text{n'est}\\
 &\quad\text{pas un segment}
 \end{aligned}
\]
(NB \deleted{\unsure{p.ex|p.}}\inserted{\unsure{c'est}} le cas si $R=\phi$ p. ex)

% slanted
\margin{faux}
Montrons dans ce cas que \deleted{\unsure{on a|est}} $\Sigma$ est \deleted{défini} \inserted{\unsure{connu|connus}} quand on connait la partition $\widehat\Sigma^*=\{A^*\mid A\in\widehat\Sigma\}$

\[
 \bigl(\text{où }\widehat\Sigma=\left\{\begin{aligned}
 &\Sigma\quad\text{si }R=\phi\\
 &\Sigma\cup\{R\}\quad\text{si }R\ne\phi
 \end{aligned}\right.\bigr)
\]
% archive sheet number: 46
